\documentclass[10pt]{article}

% Specify the margins and text size.
\setlength{\textwidth}{6.4in}
\setlength{\textheight}{9.6in}
\setlength{\oddsidemargin}{0pt}
\setlength{\evensidemargin}{0pt}
\setlength{\topmargin}{0pt}
\setlength{\hoffset}{.05in}
\setlength{\voffset}{-1in}

\setlength{\parskip}{5pt}
\setlength{\parindent}{0pt}

% Load some fonts and symbol packages
\usepackage{latexsym}
\usepackage{pifont}       % contains 'star' symbol for counterinsurgency handbook title
\usepackage{yfonts} 
\usepackage{amsmath}
\usepackage{amsfonts}

\usepackage{graphicx}     % actually, this is loaded by pstricks
\usepackage[T1]{fontenc}
\usepackage{ifthen}
\usepackage{pstricks,pst-func,pst-text,pst-node,multido,pst-plot,calc,pst-3dplot}
%\usepackage[all]{xy}
%\usepackage{animate}

% The hyperref package inserts links into the pdf.
\definecolor{MyLinkColor}{rgb}{.1,.2,1}
\definecolor{MyCiteColor}{rgb}{.1,1,.2}
\definecolor{MyURLColor}{rgb}{.4,.4,.4}
\usepackage[backref=true,pagebackref=false,hyperindex,colorlinks=true,
  linkcolor=MyLinkColor,urlcolor=MyURLColor]{hyperref}


% The tweaklist package is something I found on the web.  It provides a simple interface
% for making changes to spacing used in the itemize and enumerate environments.  Comment
% this out if you don't care to use tweaklist.
\usepackage{tweaklist}
\renewcommand{\itemhook}{\setlength{\parskip}{2pt}\setlength{\parsep}%
{1pt}\setlength{\topsep}{0pt}\setlength{\itemsep}{0pt}}

\newcommand{\N}{\ensuremath\mathbb{N}}
\newcommand{\Q}{\ensuremath\mathbb{Q}}
\newcommand{\R}{\ensuremath\mathbb{R}}
\newcommand{\C}{\ensuremath\mathbb{C}}
\newcommand{\z}{\hphantom{0}}

\begin{document}
\pagestyle{empty}

\href{http://www.shepherd.edu}{\includegraphics[height=1.75cm]{logo-high-res.eps}}
\vspace{-1.69cm}

{\small
\begin{tabular}{cl}
\hspace{5in} & Math 314:  Statistics\\
%                & Worksheet\\
                & %1 February 2012
\end{tabular}
}
\medskip

{\setlength{\baselineskip}{1.05\baselineskip}
\begin{center}
\textbf{\large  Worksheet:  Areas Under the Normal Curve}
\end{center}

\newcommand{\Z}{\hphantom{0}}
1. You can calculate the area under the normal curve and to the left of a specified value
$z$ in \texttt{R} using the command \texttt{pnorm(z)}.  Sketch the following regions then
use $\texttt{R}$ to find their areas.\vspace{-4pt}
\begin{center}
\begin{tabular}{lcl}
a) to the left of $z=0$ &\hspace{3in} & b) between $-1$ and $1$\\[.7in]
c) between $-2$ and $2$ &\hspace{3in} & d) between $-3$ and $3$\\[.7in]
e) to the left of $1.5$ &\hspace{3in} & f) to the right of $1.5$\\[.7in]
g) $-2.1<z<2$ &\hspace{3in} & h) $z>4$\\[.5in]
\end{tabular}
\end{center}
\vfill


\href{http://www.shepherd.edu}{\includegraphics[height=1.75cm]{logo-high-res.eps}}
\vspace{-1.69cm}

{\small
\begin{tabular}{cl}
\hspace{5in} & Math 314:  Statistics\\
%                & Worksheet\\
                & %1 February 2012
\end{tabular}
}
\medskip

\begin{center}
\textbf{\large  Worksheet:  Areas Under the Normal Curve}
\end{center}

1. You can calculate the area under the normal curve and to the left of a specified value
$z$ in \texttt{R} using the command \texttt{pnorm(z)}.  Sketch the following regions then
use $\texttt{R}$ to find their areas.\vspace{-4pt}
\begin{center}
\begin{tabular}{lcl}
a) to the left of $z=0$ &\hspace{3in} & b) between $-1$ and $1$\\[.7in]
c) between $-2$ and $2$ &\hspace{3in} & d) between $-3$ and $3$\\[.7in]
e) to the left of $1.5$ &\hspace{3in} & f) to the right of $1.5$\\[.7in]
g) $-2.1<z<2$ &\hspace{3in} & h) $z>4$\\
\end{tabular}
\end{center}
\vfill
\eject

2. To find the $z$ score for which $p$ percent of the area under the normal curve is
to the left of $z$, use the command \texttt{qnorm(p)}.
Find the $z$ value then sketch the region.
\begin{center}
\begin{tabular}{lcl}
a) area to the left of $z$ is 50\% &\hspace{2in} & b) area to the left of $z$ is 75\%\\[.7in]
c) area to the left of $z$ is 25\% &\hspace{2in} & d) area to the righ of $z$ is 5\%\\[.7in]
e) area to the righ of $z$ is 65\% &\hspace{2in} & f) area between $-z$ and $z$ is 25\%\\[.7in]
\end{tabular}
\end{center}

3. In 2005, men averaged about 540 on the Math SAT.  
The SD was about 120. 
Sketch a graph of the histogram of SAT scores.
Estimate the percentage of men who scored over 700.  
About what range includes 
75\% of the SAT scores were below what value?
\vspace{1.5in}


2. To find the $z$ score for which $p$ percent of the area under the normal curve is
to the left of $z$, use the command \texttt{qnorm(p)}.
Find the $z$ value then sketch the region.
\begin{center}
\begin{tabular}{lcl}
a) area to the left of $z$ is 50\% &\hspace{2in} & b) area to the left of $z$ is 75\%\\[.7in]
c) area to the left of $z$ is 25\% &\hspace{2in} & d) area to the righ of $z$ is 5\%\\[.7in]
e) area to the righ of $z$ is 65\% &\hspace{2in} & f) area between $-z$ and $z$ is 25\%\\[.7in]
\end{tabular}
\end{center}

3. In 2005, men averaged about 540 on the Math SAT.  
The SD was about 120. 
Sketch a graph of the histogram of SAT scores.
Estimate the percentage of men who scored over 700.  
About what range includes 
75\% of the SAT scores were below what value?

\end{document}
